\documentclass[11pt]{article}

\usepackage{fancyhdr}
\pagestyle{fancy}
\fancyhf{} % clear default header/footer
\fancyhead[C]{Working Notes: Claims Not Verified}
\renewcommand{\headrulewidth}{0.4pt}

\usepackage[margin=1in]{geometry}
\usepackage[T1]{fontenc}
\usepackage[utf8]{inputenc}
\usepackage{amsmath,amssymb,amsthm,mathtools}
\usepackage{booktabs}
\usepackage{enumitem}
\usepackage{xcolor}
\usepackage[colorlinks=true,linkcolor=blue!55!black,citecolor=blue!55!black,urlcolor=blue!55!black]{hyperref}
\usepackage{natbib}
\usepackage{dsfont}

\newtheorem{definition}{Definition}
\newtheorem{proposition}{Proposition}
\newtheorem{lemma}{Lemma}
\newtheorem{theorem}{Theorem}
\newtheorem{corollary}{Corollary}
\theoremstyle{remark}
\newtheorem{remark}{Remark}

\newcommand{\Acal}{\mathcal A}
\newcommand{\Scal}{\mathcal S}
\newcommand{\Zcal}{\mathcal Z}
\newcommand{\Dcal}{\mathcal D}
\newcommand{\Gcal}{\mathcal G}
\newcommand{\Ncal}{\mathcal N}
\newcommand{\E}{\mathbb E}
\newcommand{\1}{\mathds{1}}
\newcommand{\dd}{\mathrm d}

% This note: occupancy, pushforward, outcome marginal, entropy functional.
\newcommand{\occ}{d^{\pi}}
\newcommand{\pz}{p_{\pi}}
\newcommand{\Hf}[1]{H_{\mu}\!\left[#1\right]}
\newcommand{\Hl}[1]{H_{\lambda}\!\left[#1\right]}
\newcommand{\KL}{\mathrm{KL}}
\DeclareMathOperator*{\argmin}{arg\,min}
\DeclareMathOperator*{\argmax}{arg\,max}
\newcommand{\new}[1]{\textcolor{blue!70!black}{#1}}
\newcommand{\fix}[1]{\textcolor{red!75!black}{#1}}
\newcommand{\todo}[1]{\textcolor{red!70!black}{[#1]}}

\title{From Theory to Algorithm (v2): Good Regulator inspired Optimization for $\min_\pi \Hf{p_\pi}$ over Trajectories\\
\Large \textbf{Working Notes: Claims Not Verified}}
\author{Joery A. de Vries}

\begin{document}
\maketitle

\noindent
This is v1 (\texttt{theory\_to\_algorithm\_joery.tex}) rewritten against a general reference measure $\mu$ on the outcome space, with the finite alphabet and the Euclidean case recovered as instantiations in Section~\ref{sec:inst}, and with the average-reward occupancy carried alongside the discounted one.
\new{Blue} marks content that is new relative to v1; replacing sums over $\Zcal$ by integrals against $\mu$ is not marked.
Throughout, $\ln$ is the natural logarithm and every entropy is in nats; $\log$ would denote $\log_2$ and $\log_{10}$ the decimal logarithm, and neither is used.
The entropy budget and the empowerment term live in the separate note \nolinkurl{notes_joery/entropy_budget_empowerment.tex} and are unaffected by anything here.

\section{Setting}
\label{sec:setting}

The companion draft on action-sufficient representation suggests a minimization of the entropy for a regulated outcome $Z$ for a \emph{single} decision: $p(s)$ is given, a policy $\pi(a\mid s)$ is chosen, and $\pz=\sum_{s,a}p(s)\pi(a\mid s)K(\cdot\mid s,a)$.
To turn this into a functional reinforcement learning algorithm let's consider control over trajectories where the state distribution is induced marginally by the policy.

\paragraph{Two occupancies.}
Consider an MDP $(\Scal,\Acal,P,\rho_0,\gamma)$ with transition probabilities $P(s'\mid s,a)$, initial distribution $\rho_0(s)$, discount $\gamma\in[0,1]$, and a stationary policy $\pi(a\mid s)$.
We keep $\Scal$ and $\Acal$ finite, since the algorithm below is tabular in them; nothing in this section uses it.
For a stationary policy write $\Pr_\pi(s_t=s)$ for the law of the state at time $t$ from $s_1\sim\rho_0$.
The two state--action occupancies are
\begin{equation}
  \occ(s,a)=\lim_{T\to\infty}\frac1T\sum_{t=1}^{T}\Pr_\pi(s_t=s)\,\pi(a\mid s),
  \qquad
  \occ_\gamma(s,a)=(1-\gamma)\sum_{t\ge1}\gamma^{t-1}\Pr_\pi(s_t=s)\,\pi(a\mid s),
  \label{eq:occ}
\end{equation}
the long-run and the discounted fraction of time spent in each pair.
{\color{blue!70!black}
We take over the setting of \citet{Zahavy-reward21}: an infinite-horizon finite MDP, with the average and the discounted case distinguished only through the polytopes of admissible occupancies,
\begin{equation}
\begin{aligned}
  \Dcal&=\Big\{d\ge0:\ \textstyle\sum_{s,a}d(s,a)=1,\ \ \sum_a d(s',a)=\sum_{s,a}d(s,a)\,P(s'\mid s,a)\ \ \forall s'\Big\},\\
  \Dcal_\gamma&=\Big\{d\ge0:\ \textstyle\sum_a d(s',a)=(1-\gamma)\rho_0(s')+\gamma\sum_{s,a}d(s,a)\,P(s'\mid s,a)\ \ \forall s'\Big\},
\end{aligned}
\label{eq:polytope}
\end{equation}
their Definition~1, after \citet[Ch.~6 and~8]{Puterman-markov94}.
Each is a polytope, hence convex and compact, its vertices are the occupancies of deterministic policies, and a policy is read off an occupancy as $\pi(a\mid s)=d(s,a)/\sum_{a'}d(s,a')$ wherever the denominator is positive.
The constraint defining $\Dcal$ is that of $\Dcal_\gamma$ at $\gamma=1$ with the initial distribution dropping out, and $\occ_\gamma\to\occ$ as $\gamma\to1$, since an Abel mean and a Ces\`aro mean of the same bounded sequence agree whenever the latter exists.\footnote{In a finite MDP the Ces\`aro limit in~\eqref{eq:occ} exists for every stationary policy. That it does not depend on $\rho_0$, and that every point of $\Dcal$ is attained by some stationary policy, is where a recurrence assumption on $P^\pi$ (unichain) enters; like \citet{Zahavy-reward21} we work with the polytope directly and make no further distinction between the two settings.}
The unsubscripted $\occ$ and $\Dcal$ are the average-case objects; the discounted ones carry their $\gamma$, and where a statement holds for both we write $\occ$, $\Dcal$ and say so.
Every statement below uses only that $\Dcal$ is a polytope and that $\occ\mapsto\occ K$ is linear.
}
A reward $r(s,a)\in[0,1]$ has the return
\begin{equation}
  R(\pi)=\sum_{s,a}\occ(s,a)\,r(s,a),
  \label{eq:return}
\end{equation}
which is the average reward $\lim_T\frac1T\E\sum_{t\le T}r_t$ in one case and $(1-\gamma)$ times the discounted return in the other (their Proposition~1), and is \emph{linear} in $\occ$ in both.
Standard RL maximises it; we write $R^\ast=\max_\pi R(\pi)$ for the maximum and $R$ for $R(\pi)$ when the policy is clear from context.
Policy search is thus a linear programme over $\Dcal$, in practice run as gradient ascent on a parametric policy $\pi_\theta$ or through dynamic programming.

\paragraph{Outcomes and the reference measure.}
Let every step emit an outcome $z\in\Zcal$ through the channel $K(\dd z\mid s,a)$ of the companion drafts.
{\color{blue!70!black}
Fix a $\sigma$-finite reference measure $\mu$ on $\Zcal$ and assume every $K(\cdot\mid s,a)$ has a density with respect to it, which we denote by the same letter, $K(\dd z\mid s,a)=K(z\mid s,a)\,\mu(\dd z)$.
The outcome marginal is the composition of the occupancy with the kernel,
\begin{equation}
  \pz(z)=(\occ K)(z)=\sum_{s,a}\occ(s,a)\,K(z\mid s,a),
  \label{eq:push}
\end{equation}
a density with respect to $\mu$. When the channel is a deterministic map $\zeta$, $\pz=\zeta_{\#}\occ$ is a pushforward of measures; for a general kernel it is the composition of a measure with a kernel, and we write $\occ K$ throughout.
Entropy is taken against $\mu$,
\begin{equation}
  \Hf{p}=-\int_{\Zcal} p\ln p\,\dd\mu ,
  \label{eq:ent}
\end{equation}
which is Shannon entropy for counting $\mu$ and differential entropy for Lebesgue $\mu$.
Differential entropy depends on $\mu$, so the choice of reference is part of the objective; Section~\ref{sec:inst} makes that dependence explicit and uses it to put a cost into the problem without touching the channel.
}

\paragraph{The objective.}
Conant and Ashby's criterion~\citep{Conant-good70} becomes
\begin{equation}
  \min_{\pi}\Hf{\pz}=\min_{\occ\in\Dcal}\Hf{\occ K}.
  \label{eq:obj}
\end{equation}
The map $\occ\mapsto\occ K$ is linear and $p\mapsto-p\ln p$ is concave pointwise, so~\eqref{eq:obj} minimises a concave functional over a convex set.
{\color{blue!70!black}
Before minimising we need the objective to be finite, and for a non-atomic $\mu$ that is not automatic: differential entropy is unbounded below, a density that piles its mass onto a set of vanishing $\mu$-volume has entropy tending to $-\infty$, and a channel that puts mass on a $\mu$-null set, a single point of $\mathbb R^d$ say, has no density at all.
What keeps the objective finite is that the policy can only mix the channel densities.
By~\eqref{eq:push} the marginal is a convex combination of the $|\Scal||\Acal|$ channel densities with weights $\occ$, so by concavity
\begin{equation}
  \Hf{\occ K}\ \ge\ \sum_{s,a}\occ(s,a)\,\Hf{K(\cdot\mid s,a)}\ \ge\ \min_{s,a}\Hf{K(\cdot\mid s,a)} ,
  \label{eq:floor}
\end{equation}
the middle term being $H(Z\mid S,A)$ under the joint law $\occ\otimes K$.
The objective is thus finite exactly when every channel density has finite entropy, that is, when no state--action pair can concentrate the outcome on zero $\mu$-volume.
For counting $\mu$ every non-empty set has volume at least one and this holds trivially.
For Lebesgue $\mu$ on $\mathbb R^d$ it excludes deterministic channels, since a point outcome has zero volume.
The requirement is on the pair $(K,\mu)$ and not on the world: coarse-graining the outcome, mapping a whole interval to one symbol, gives $\mu$ an atom and readmits a deterministic channel into it.
Either way $\min_{s,a}\Hf{K(\cdot\mid s,a)}$ is the floor the regulator cannot go below.
}
Given a finite value, the minimum of a concave function over a polytope sits at a vertex, since any point is a convex combination of vertices and concavity gives a value no smaller than the smallest vertex value.
So a deterministic stationary policy is optimal and the good-regulator conclusion $H(A\mid S)=0$ carries over unchanged.
The objective is not linear in $\occ$, however, so it is not the return of any reward $r(s,a)$ and local optima are expected.
Behaviourally, the policy should reach a set of states whose outcomes are indistinguishable and low-noise, and stay there.

\paragraph{Occupancy homeostasis.}
Consider for a moment that the channel $K$ is the identity, i.e., $\pz\equiv\occ$ with counting $\mu$.
This is the objective of SMiRL~\citep{Berseth-smirl21}, a homeostasis objective over the state--action occupancy.
Such an objective typically makes the policy stay still unless extraneous factors (environment, other agents, kernel noise) inject entropy, and in practice ``doing nothing'' or reaching a terminal state as fast as possible can be entropy-optimal.
Suppose $\rho_0$ is a point mass and $P$ and $K$ are deterministic; then no randomness enters the trajectory, a policy that reaches an absorbing pair and holds it pays entropy only for the transient, and the objective says nothing about which state is worth holding.
What gives the objective meaning is a persistent disturbance, the disturbance $D$ that regulation exists to counteract in~\citet{Conant-good70}, and we assume $P$ or $K$ carries one.
The working assumption is then that homeostasis can be formulated to favour a desirable set of outcomes.

\paragraph{Reward as a deterministic channel.}
Entropy is invariant to relabelling outcomes, so no assignment of symbols can make one outcome preferable to another by name; what a channel controls is how much of $\mu$ each state--action pair spreads its mass over.
A deterministic channel $z=\zeta(s,a)$ has $K(\cdot\mid s,a)=\delta_{\zeta(s,a)}$, \new{which by~\eqref{eq:floor} is admissible only where $\mu$ has an atom at $\zeta(s,a)$, so take counting $\mu$ here.}
Then $\pz=\zeta_{\#}\occ$ is a pushforward of measures, \new{the one place in this note where the word is literal,} and $\Hf{\pz}$ is never linear in $\occ$, so a linear reward cannot be recovered this way.
What $\zeta$ does control is the granularity of the partition $\zeta^{-1}(z)$ of $\Scal\times\Acal$, which is where desirability can enter.
Take $\Gcal=\{z_0\}$, let $\zeta$ map every desirable pair to $z_0$ and every undesirable pair to a private symbol, and assume a disturbance in $P$ that moves the agent between nearby pairs.
Write $R=\sum_{(s,a)\in\zeta^{-1}(z_0)}\occ(s,a)$ for the mass on the desirable pairs, the many pairs that share the symbol summed into one number and the return~\eqref{eq:return} of the indicator reward $r=\1\{\zeta(s,a)=z_0\}$, and $\tilde d=\occ/(1-R)$ for the normalised occupancy on the rest, so that $\pz(z_0)=R$ and $\pz(\zeta(s,a))=(1-R)\,\tilde d(s,a)$ for every other pair.
Substituting into~\eqref{eq:ent} and splitting the logarithm of the product,
\begin{equation}
  \Hf{\pz}
  =-R\ln R-\sum_{(s,a)\notin\zeta^{-1}(z_0)}(1-R)\,\tilde d(s,a)\ln\big((1-R)\,\tilde d(s,a)\big)
  =h_b(R)+(1-R)\,\Hf{\tilde d},
  \label{eq:det-channel-entropy}
\end{equation}
using $\sum\tilde d=1$, with $h_b(R)=-R\ln R-(1-R)\ln(1-R)$ the binary entropy.
The first term is at most $\ln2$; the cost sits in the second, the entropy of how the disturbance spreads the remaining mass over individually labelled pairs.
A policy holding the desirable pairs is moved among pairs that share $z_0$ at no entropy cost, a policy holding an undesirable pair is moved among distinct symbols and pays, so holding the desirable set is preferred.
This is an ordinal preference that depends on the disturbance being persistent and the desirable set being large enough that the disturbance rarely leaves it, and it cannot express a continuous reward.

\paragraph{Reward as a stochastic channel.}
Alternatively the disturbance can be placed inside the channel, which \new{also lifts the restriction to atomic $\mu$}.
{\color{blue!70!black}
Let the desirable set be a measurable $\Gcal\subseteq\Zcal$ with $0<\mu(\Gcal)<\infty$ and $0<\mu(\Gcal^c)<\infty$, so $\mu$ is finite.
For a scalar $\rho\in[0,1]$ define the density
\begin{equation}
  K_\rho(z)=\rho\,\frac{\1\{z\in \Gcal\}}{\mu(\Gcal)}+(1-\rho)\,\frac{\1\{z\in \Gcal^c\}}{\mu(\Gcal^c)} ,
  \label{eq:K-family}
\end{equation}
mass $\rho$ spread evenly over $\Gcal$ and the rest spread evenly over the complement.
Its entropy is computed on the two pieces, on each of which the density is constant,
\begin{align}
  \Hf{K_\rho}
  &=-\int_\Gcal\frac{\rho}{\mu(\Gcal)}\ln\frac{\rho}{\mu(\Gcal)}\,\dd\mu
    -\int_{\Gcal^c}\frac{1-\rho}{\mu(\Gcal^c)}\ln\frac{1-\rho}{\mu(\Gcal^c)}\,\dd\mu \notag\\
  &=-\rho\ln\frac{\rho}{\mu(\Gcal)}-(1-\rho)\ln\frac{1-\rho}{\mu(\Gcal^c)} \notag\\
  &=h_b(\rho)+\rho\ln\mu(\Gcal)+(1-\rho)\ln\mu(\Gcal^c)=:F(\rho),
  \label{eq:F}
\end{align}
with $h_b$ the binary entropy as in~\eqref{eq:det-channel-entropy}.
Read~\eqref{eq:F} as a two-stage draw: a coin with bias $\rho$ decides whether the outcome lands in $\Gcal$, at cost $h_b(\rho)$, and a position inside the chosen set is then drawn uniformly, at cost the natural logarithm of its volume, so $F(\rho)=h_b(\rho)+\E[\ln\text{volume of the chosen set}]$.
In v1 the desirable set was a single atom, $\ln\mu(\Gcal)=\ln1=0$, and the position cost inside $\Gcal$ was invisible, leaving $h_b(R)+(1-R)\ln N$; the term $\rho\ln\mu(\Gcal)$ is the cost of not knowing where inside $\Gcal$ the outcome sits, which a point hides.
Adding and subtracting $\ln\mu(\Zcal)=\ln\big(\mu(\Gcal)+\mu(\Gcal^c)\big)$ and collecting the $\rho$ and $1-\rho$ terms,
\begin{align}
  F(\rho)
  &=\ln\mu(\Zcal)-\Big[\rho\ln\rho+(1-\rho)\ln(1-\rho)+\ln\mu(\Zcal)-\rho\ln\mu(\Gcal)-(1-\rho)\ln\mu(\Gcal^c)\Big] \notag\\
  &=\ln\mu(\Zcal)-\Big[\rho\ln\frac{\rho}{\mu(\Gcal)/\mu(\Zcal)}+(1-\rho)\ln\frac{1-\rho}{\mu(\Gcal^c)/\mu(\Zcal)}\Big] \notag\\
  &=\ln\mu(\Zcal)-\KL\Big(\mathrm{Bern}(\rho)\,\Big\|\,\mathrm{Bern}\big(\mu(\Gcal)/\mu(\Zcal)\big)\Big),
  \label{eq:F-kl}
\end{align}
the bracket being the binary $\KL$ between the hit rate $\rho$ and the volume fraction $\mu(\Gcal)/\mu(\Zcal)$, since $\mu(\Gcal^c)/\mu(\Zcal)=1-\mu(\Gcal)/\mu(\Zcal)$.
So a member of the family has the entropy of the uniform law on $\Zcal$ minus how far its hit rate is from the volume fraction of $\Gcal$: $F$ is maximal, and the member uniform on $\Zcal$, at $\rho=\mu(\Gcal)/\mu(\Zcal)$, and minimising $F$ pushes the hit rate away from the volume fraction in either direction, which is the concavity picture below in one line.
Both pieces of~\eqref{eq:K-family} are affine in $\rho$, so for weights $\nu_i$ on the simplex and $\bar\rho=\sum_i\nu_i\rho_i$,
\begin{equation}
  \sum_i \nu_i K_{\rho_i}(z)
  =\bar\rho\,\frac{\1\{z\in \Gcal\}}{\mu(\Gcal)}+(1-\bar\rho)\,\frac{\1\{z\in \Gcal^c\}}{\mu(\Gcal^c)}
  =K_{\bar\rho}(z),
  \label{eq:affine}
\end{equation}
mixing the parameters is the same as mixing the densities.
Now let the reward $r(s,a)$ set the channel, $K(\cdot\mid s,a)=K_{r(s,a)}$.
Applying~\eqref{eq:affine} with $\nu=\occ$ in~\eqref{eq:push} gives $\pz=K_{R}$ with $R=R(\pi)$ the return~\eqref{eq:return}, so the marginal depends on the occupancy only through $R$, and by~\eqref{eq:F} the homeostasis objective factors through the return,
\begin{equation}
  \min_{\pi}\Hf{\occ K}=\min_{\pi}F\big(R(\pi)\big).
  \label{eq:factor}
\end{equation}
$F$ is strictly concave, being $h_b$ plus a term affine in $\rho$, so on $[0,R^\ast]$ it lies above the smaller endpoint value, $F(R)\ge\min\{F(0),F(R^\ast)\}$, with equality only at the endpoints.
The two endpoint values are $F(0)=\ln\mu(\Gcal^c)$ and $F(1)=\ln\mu(\Gcal)$: the entropy of holding a set is the natural logarithm of its size, and entropy sees $\Gcal$ only through its size.
Hence entropy minimisers are return maximisers exactly when
\begin{equation}
  F(R^\ast)<F(0)
  \quad\Longleftrightarrow\quad
  h_b(R^\ast)<R^\ast\ln\frac{\mu(\Gcal^c)}{\mu(\Gcal)} ,
  \label{eq:condition}
\end{equation}
obtained by subtracting $\ln\mu(\Gcal^c)$ from both sides of $h_b(R^\ast)+R^\ast\ln\mu(\Gcal)+(1-R^\ast)\ln\mu(\Gcal^c)<\ln\mu(\Gcal^c)$.
The condition is on the volume ratio of the undesirable to the desirable set, in units of $\mu$.
Since $h_b\le\ln2$ it holds whenever $\mu(\Gcal^c)/\mu(\Gcal)>2^{1/R^\ast}$, and it fails whenever $\Gcal$ is the larger set, in which case ``reliable failure'' is the entropy minimiser.
Subtracting $\ln\mu(\Gcal)$ from $F$ and dividing by the logarithm of the volume ratio,
\begin{equation}
  \frac{F(R)-\ln\mu(\Gcal)}{\ln\big(\mu(\Gcal^c)/\mu(\Gcal)\big)}
  =1-R+\frac{h_b(R)}{\ln\big(\mu(\Gcal^c)/\mu(\Gcal)\big)}
  \ \longrightarrow\ 1-R
  \label{eq:limit}
\end{equation}
as the ratio grows, so the normalised objective converges to $1-R(\pi)$ and minimising it recovers the standard LP.
Where in $\Gcal$ the mass sits is invisible to~\eqref{eq:F}; it is decided by how $r(s,a)$ depends on the state, which Section~\ref{sec:inst} makes concrete for a disturbance that leaks mass out of $\Gcal$.
}


\section{Practical Algorithm v2}
\label{sec:algorithm}

\paragraph{The policy needs more than the state.}
Section~\ref{sec:setting} showed that a homeostasis-optimal policy can be deterministic and stationary, so $\pi(a\mid s)$ suffices to represent the optimum.
\textbf{It does not suffice to find it.}
For a general channel the objective is not a fixed MDP reward, so there is no single MDP to solve; this is the insight of convex RL~\citep{Hazan-provably19,Zahavy-reward21}, from which we borrow.
The algorithm below is a policy iteration in which each step solves an MDP whose reward depends on the outcome marginal of the current policy, so we write $\pi(a\mid s,w)$ with $w$ a summary of $\pz$.
This is sufficient because $\occ$ enters the objective only through $\pz=\occ K$.

\paragraph{Policy iteration through the Gibbs form.}
For any density $q$ with respect to $\mu$,
\begin{equation}
  \E_{z\sim p}\big[-\ln q(z)\big]=\Hf{p}+\KL(p\,\|\,q)\ \ge\ \Hf{p},
  \label{eq:gibbs}
\end{equation}
with equality iff $q=p$, so $\Hf{p}=\min_q\E_{p}[-\ln q(z)]$; note that the reference cancels in the divergence, so the right-hand side does not depend on $\mu$ once $\Hf{p}$ does.
Apply this to $p=\pz$ and expand the expectation through the channel,
\begin{align}
  \E_{z\sim\pz}\big[-\ln q(z)\big]
  &=\sum_{s,a}\occ(s,a)\int K(z\mid s,a)\big(-\ln q(z)\big)\,\mu(\dd z)=-\sum_{s,a}\occ(s,a)\,r_q(s,a),
  \label{eq:gibbs-reward}\\[2pt]
  r_q(s,a)
  &:=\int K(z\mid s,a)\ln q(z)\,\mu(\dd z)=\E_{z\sim K(\cdot\mid s,a)}\big[\ln q(z)\big] . \notag
\end{align}
Hence
\begin{equation}
  \min_{\pi}\Hf{\occ K}=\min_{\pi}\ \min_{q}\ \Big(-\sum_{s,a}\occ(s,a)\,r_q(s,a)\Big),
  \label{eq:min-min}
\end{equation}
and for fixed $q$ the inner problem is an ordinary MDP with reward $r_q$, linear in $\occ$, in the discounted or the average-reward sense according to which occupancy is in use.
Alternating the two minimisations gives the iteration: given $\pi_k$, the inner minimiser is $q_k=p_{\pi_k}$, which we estimate by $\hat q_k$ from observed outcomes, and we \emph{improve} $\pi_k$ on the MDP with reward $r_{\hat q_k}$ to obtain $\pi_{k+1}$.
The iteration is monotone,
\begin{equation}
  \Hf{p_{\pi_{k+1}}}
  \ \le\ \E_{p_{\pi_{k+1}}}\big[-\ln q_k\big]
  \ \le\ \E_{p_{\pi_k}}\big[-\ln q_k\big]
  \ =\ \Hf{p_{\pi_k}},
  \label{eq:mm-descent}
\end{equation}
by~\eqref{eq:gibbs}, then improvement under the frozen reward, then $q_k=p_{\pi_k}$; with $\hat q_k$ in place of $q_k$ the last equality picks up $\KL(p_{\pi_k}\,\|\,\hat q_k)$ and the descent holds up to that error.
The objective is concave, so the iteration converges to a local optimum determined largely by $\pi_0$, restarts are needed for anything global, and minimising a concave function over a polytope is NP-hard in general~\citep{Pardalos-quadratic91}.

\paragraph{\new{The average-reward inner problem.}}
{\color{blue!70!black}
With the average occupancy $\occ$ the inner MDP is an average-reward MDP.
Its gain is $g=\sum_{s,a}\occ(s,a)r_q(s,a)$, which by~\eqref{eq:gibbs-reward} is $-\E_{\pz}[-\ln q]$, minus the cross-entropy, and improvement is judged by the differential action value
\begin{equation}
  Q^\pi(s,a)=r_q(s,a)-g+\sum_{s'}P(s'\mid s,a)\sum_{a'}\pi(a'\mid s')\,Q^\pi(s',a'),
  \label{eq:avg-q}
\end{equation}
defined up to an additive constant that the softmax actor below absorbs.
A policy that improves on $\pi$ under $Q^\pi$ at every state has no smaller gain under the recurrence assumption of the footnote in Section~\ref{sec:setting}~\citep[\S8.6]{Puterman-markov94}; \citet{Zahavy-reward21} instead treat the average-reward solver as a black-box best response, and either reading gives the improvement that~\eqref{eq:mm-descent} asks for.
On the estimation side nothing changes except the weights: the outcome estimator of v1 (Proposition~1 there) at $\gamma=1$ and forgetting $\lambda=0$ is the plain sample mean $\beta_t=1/(t+1)$, and it estimates $\occ K$.
For an episodic task run with restarts, $\occ$ is the occupancy of the chain that includes the restart transition, so the per-episode average is already the right object and no discount bookkeeping is needed.
}

\paragraph{\new{Representing $q$: the task vector is a function.}}
{\color{blue!70!black}
In v1 the task vector was $w=\ln\hat q\in\mathbb R^{|\Zcal|}$.
For general $\mu$ the object $\ln\hat q$ is a function on $\Zcal$, and a finite parametrisation is needed.
Take an exponential family with respect to $\mu$ with features $\phi:\Zcal\to\mathbb R^m$,
\begin{equation}
  \hat q_w(z)=\exp\big(w^\top\phi(z)-\Lambda(w)\big),
  \qquad
  \Lambda(w)=\ln\int e^{w^\top\phi}\,\dd\mu ,
  \label{eq:expfam}
\end{equation}
so that $\ln\hat q_w$ is affine in $\phi$ and the Gibbs reward~\eqref{eq:gibbs-reward} becomes
\begin{equation}
  r_q(s,a)=w^\top\E_{K(\cdot\mid s,a)}\big[\phi(Z)\big]-\Lambda(w),
  \label{eq:expfam-reward}
\end{equation}
a linear function of the channel's feature expectation plus a constant.
Its discounted action value is then a successor feature in $\phi$,
\begin{equation}
  Q_w(s,a)=w^\top\psi^\pi(s,a)-\Lambda(w)\,m^\pi(s,a),
  \qquad
  \psi^\pi=\E_{K}[\phi]+\gamma P^\pi\psi^\pi,
  \qquad
  m^\pi=1+\gamma P^\pi m^\pi ,
  \label{eq:sf-q}
\end{equation}
where $m^\pi$ is the remaining discounted horizon, equal to $1/(1-\gamma)$ in the continuing case and to the remaining episode mass otherwise.
If $\phi$ contains the constant feature $1$ then $m^\pi$ is one coordinate of $\psi^\pi$ and~\eqref{eq:sf-q} is invariant under $w\mapsto w+c\,e_1$: the constant term shifts $w^\top\psi^\pi$ by $c\,m^\pi$ and $\Lambda$ by $c$, and the two cancel, at every stage of learning and not only at convergence.
This is the exact form of the baseline correction $\bar w$ of v1, which subtracted the mean of $\ln\hat q$ but not the matching $\Lambda$ term.
In the average-reward case the constant drops out of $r_q-g$ altogether, and~\eqref{eq:sf-q} becomes $Q_w=w^\top\psi^\pi$ with the centred recursion $\psi^\pi=\E_K[\phi]-\E_{\pz}[\phi]+P^\pi\psi^\pi$.
Fitting $\hat q_w$ to outcomes is maximum likelihood in the family, which matches moments, $\nabla\Lambda(w)=\E_{\hat q_w}[\phi]=\bar\phi$, with $\bar\phi$ the $\omega$-weighted sample mean of $\phi(z_t)$ produced by the estimator of v1 with $\phi(z_t)$ in place of the indicator $e_{z_t}$; its proof never used that $e_{z_t}$ is an indicator.
One-hot features on a finite alphabet make $\bar\phi$ the histogram and $w=\ln\hat q$ up to the gauge, recovering v1.
Bin indicators on a continuous $\Zcal$ make it a histogram density, $\phi(z)=(z,z^2)$ makes it a Gaussian fit.
}

\paragraph{Feature instantiation.}
Tables are logits $\theta_t\in\mathbb R^{|\Scal|\times|\Acal|}$ with $\pi_t=\mathrm{softmax}(\theta_t)$, successor features $\psi_t\in\mathbb R^{|\Scal|\times|\Acal|\times m}$, a moment estimate $\bar\phi_t\in\mathbb R^m$ with discount weight $\omega_t$ and accumulator $W_t$ as in v1, and the model parameter $w_t$ solving $\nabla\Lambda(w_t)=\bar\phi_t$.
At transition $t$ the environment emits $z_t\sim K(\cdot\mid s_t,a_t)$, $d_t\in\{0,1\}$ marks the start of an episode, $\omega_t=1$ if $d_t=1$ and $\gamma\omega_{t-1}$ otherwise, and $a_{t+1}\sim\pi_t$.
With learning rates $\alpha,\eta$ and episodic forgetting $\lambda\in[0,1)$,
\begin{align}
  \psi_{t+1}(s_t,a_t) &= \psi_t(s_t,a_t)+\alpha\big[\phi(z_t)+\gamma\,(1-d_{t+1})\,\psi_t(s_{t+1},a_{t+1})-\psi_t(s_t,a_t)\big],
  \label{eq:tab-sf}\\
  \bar\phi_{t+1} &= \bar\phi_t+\beta_t\big(\phi(z_t)-\bar\phi_t\big),
  \qquad \beta_t=\frac{\omega_t}{W_{t+1}},
  \qquad W_{t+1}=(1-\lambda\,d_t)\,W_t+\omega_t,
  \label{eq:tab-w}\\
  \theta_{t+1}(s_t,\cdot) &= \theta_t(s_t,\cdot)+\eta\,Q_{w_t}(s_t,\cdot),
  \qquad
  Q_{w_t}(s,a)=w_t^\top\psi_t(s,a)-\Lambda(w_t)\,m_t(s,a) ,
  \label{eq:tab-pi}
\end{align}
with $m_t$ the constant-feature coordinate of $\psi_t$, and with the centred recursion and no $\Lambda$ term at $\gamma=1$.
For a fixed $w$, \eqref{eq:tab-sf} is SARSA on the scalar reward $w^\top\phi(z_t)$, and the actor is the natural policy gradient for a tabular softmax policy~\citep{Kakade-natural01,Agarwal-theory21}, the multiplicative update $\pi_{t+1}(\cdot\mid s)\propto\pi_t(\cdot\mid s)\exp(\eta\,Q_{w_t}(s,\cdot))$, which improves $\sum_a\pi(a\mid s)Q_{w_t}(s,a)$ at $s_t$ and keeps full support.
Both $\psi_t$ and $\bar\phi_t$ estimate quantities of the current policy, so $\eta$ should be small relative to $\alpha$ and $\lambda$; a stale value or a stale model appears as the $\KL$ term in~\eqref{eq:mm-descent}.
Within one episode the weights decay by $\gamma$ per step and forgetting fires only at episode boundaries; at $\gamma=1$, $\lambda=0$ the weights are flat and $\beta_t=1/(t+1)$, the average-reward estimator of the previous paragraph.

\paragraph{\new{The reward channel is a one-parameter family.}}
{\color{blue!70!black}
For the channel~\eqref{eq:K-family} the whole construction collapses to a scalar.
The family is the exponential family with the single feature $\phi(z)=\1\{z\in \Gcal\}$, since for a scalar $w$
\begin{equation}
  e^{w\1\{z\in \Gcal\}}\big/\big(\mu(\Gcal)e^{w}+\mu(\Gcal^c)\big)=K_\rho(z)
  \quad\text{with}\quad
  \rho=\frac{\mu(\Gcal)e^{w}}{\mu(\Gcal)e^{w}+\mu(\Gcal^c)},
  \qquad
  w=\ln\frac{\rho}{1-\rho}+\ln\frac{\mu(\Gcal^c)}{\mu(\Gcal)} .
  \label{eq:family-expfam}
\end{equation}
The model parameter is the log-odds of the estimated return, shifted by the logarithm of the volume ratio, and $\E_{K(\cdot\mid s,a)}[\phi]=r(s,a)$.
Writing $\ln K_\rho$ out from~\eqref{eq:K-family} with $\1\{z\in \Gcal^c\}=1-\1\{z\in \Gcal\}$,
\begin{equation}
  \ln K_\rho(z)=\1\{z\in \Gcal\}\Big(\ln\frac{\rho}{\mu(\Gcal)}-\ln\frac{1-\rho}{\mu(\Gcal^c)}\Big)+\ln\frac{1-\rho}{\mu(\Gcal^c)}
  \quad\Longrightarrow\quad
  r_q(s,a)=w\,r(s,a)+\ln\frac{1-\rho}{\mu(\Gcal^c)} ,
  \label{eq:family-reward}
\end{equation}
which is~\eqref{eq:expfam-reward} with $\Lambda(w)=\ln\big(\mu(\Gcal^c)/(1-\rho)\big)$.
So on the reward channel the inner MDP of every iteration is the original reward $r$ scaled by the scalar $w_k$ plus a constant, and the algorithm is SARSA on $r$ with a time-varying gain and the softmax actor.
The gain is positive, and the iteration seeks $\Gcal$, exactly when $\rho_k>\mu(\Gcal)/\mu(\Zcal)$, by rearranging $w>0$ in~\eqref{eq:family-expfam}, which is the point where $F$ is maximal in~\eqref{eq:F-kl}: the estimated return has to exceed the volume fraction of the desirable set.
Below that fraction the iteration moves away from $\Gcal$, and $R=0$ is a fixed point, which is what the $\varepsilon$-floor on $\hat q$ in v1 perturbs.
This local condition is weaker than the global one~\eqref{eq:condition}, as it must be for a majorisation--minimisation scheme that converges to local optima.
In the counting case it reads $\rho_k>1/(N+1)$.
}


\section{Instantiations}
\label{sec:inst}

\paragraph{Counting measure on a finite alphabet.}
Let $\mu$ be counting measure on $\Zcal=\{z_0\}\cup\Ncal$ with $|\Ncal|=N$, and $\Gcal=\{z_0\}$.
Then $\mu(\Gcal)=1$, $\mu(\Gcal^c)=N$, and the objects of this note reduce to those of v1 line by line:

\begin{center}
\begin{tabular}{lll}
\toprule
here & becomes & in v1 \\
\midrule
$K_\rho$, eq.~\eqref{eq:K-family} & $\rho\,\1\{z=z_0\}+(1-\rho)\tfrac1N\1\{z\in\Ncal\}$ & eq.~(8), the stochastic channel \\
$F(\rho)$, eq.~\eqref{eq:F} & $h_b(\rho)+(1-\rho)\ln N$ & eq.~(10) \\
condition~\eqref{eq:condition} & $h_b(R^\ast)<R^\ast\ln N$ & eq.~(12) \\
limit~\eqref{eq:limit} & $F(R)/\ln N\to1-R$ & eq.~(14) \\
$\phi$, $\bar\phi_t$ & one-hot $e_z$, histogram $\hat q_t$ & eq.~(20) and Prop.~1 \\
$Q_w$, eq.~\eqref{eq:sf-q} & $\psi^\top\ln\hat q$, gauge-fixed by $\Lambda$ & eq.~(22) with $\bar w$ \\
gain $w_k$, eq.~\eqref{eq:family-expfam} & $\ln\hat q(z_0)-\ln\big((1-\hat q(z_0))/N\big)$ & implicit in $\psi^\top w$ \\
\bottomrule
\end{tabular}
\end{center}

\noindent
The last row is worth seeing once: with $\psi(s,a,z_0)=Q^\pi_r(s,a)$ and each noise coordinate $\big(m^\pi-Q^\pi_r\big)/N$, v1's $\psi^\top w$ equals $Q^\pi_r\big(\ln\hat q(z_0)-\ln\tfrac{1-\hat q(z_0)}{N}\big)+m^\pi\ln\tfrac{1-\hat q(z_0)}{N}$, which is $w_kQ^\pi_r-\Lambda(w_k)m^\pi$.
DeepSea as in v1: episodes of length $L$, treasure on the last step, $R^\ast=(1-\gamma)\gamma^{L-1}/(1-\gamma^L)$, and the treasure-seeking regime needs $N>\exp\big(h_b(R^\ast)/R^\ast\big)$.

\paragraph{\new{Lebesgue measure on a bounded box.}}
{\color{blue!70!black}
Let $\Zcal\subset\mathbb R^d$ be a bounded box with Lebesgue measure $\lambda$ as reference, so that $\Hf{\cdot}$ is differential entropy, and let $\Gcal$ be an interval $[z_0-\delta,z_0+\delta]$ in $d=1$, so $\lambda(\Gcal)=2\delta$ and $\lambda(\Gcal^c)=\lambda(\Zcal)-2\delta$.
Condition~\eqref{eq:condition} reads $h_b(R^\ast)<R^\ast\ln\big((\lambda(\Zcal)-2\delta)/2\delta\big)$, the alphabet size of v1 replaced by a volume ratio, and it can only hold if $\Gcal$ occupies less than half the box.
For the disturbance take an additive noise $\epsilon$ with a symmetric, unimodal density $f$ around a set-point $\zeta(s,a)$ that the state--action pair commands, and let the channel be
\begin{equation}
  \tilde z=\zeta(s,a)+\epsilon,
  \qquad
  Z\sim\mathrm{Unif}(\Gcal)\ \text{if }\tilde z\in \Gcal,
  \qquad
  Z\sim\mathrm{Unif}(\Gcal^c)\ \text{otherwise},
  \label{eq:leak-channel}
\end{equation}
which is~\eqref{eq:K-family} with $r(s,a)=\Pr\big(\zeta(s,a)+\epsilon\in \Gcal\big)$, the probability that the disturbance does not push the outcome out of the band.
Everything in Sections~\ref{sec:setting} and~\ref{sec:algorithm} applies with this $r$, and the entropy minimiser maximises $R$, that is, minimises the leak.
Which set-point that selects is a one-line computation: for $x=\zeta(s,a)$,
\begin{equation}
  r(x)=\int_{z_0-\delta-x}^{z_0+\delta-x}f(u)\,\dd u,
  \qquad
  r'(x)=f(z_0-\delta-x)-f(z_0+\delta-x),
  \label{eq:leak-derivative}
\end{equation}
which vanishes at $x=z_0$ by symmetry and is negative for $x>z_0$ by unimodality, since $|z_0-\delta-x|>|z_0+\delta-x|$ there.
The regulator therefore holds the centre of the band, not merely the band, and the strength of that preference is set by the noise width relative to $\delta$.
This is target-seeking derived from homeostasis rather than imposed on it.
Two remarks keep the derivation honest.
If the raw $\tilde z$ were emitted instead, with $\Zcal=\mathbb R$, then $\pz$ would be the law of $\zeta$ convolved with $f$, differential entropy is translation invariant, and no set-point would be preferred: it is the redraw on leaving $\Gcal$ that makes position in $\Gcal$ matter.
And the uniform redraw inside $\Gcal$ is what keeps the channel inside the family; a channel that concentrates inside $\Gcal$ lowers $\Hf{Z\mid S,A}$ further, but by~\eqref{eq:floor} not below the entropy of $f$.
For the algorithm the single feature $\phi(z)=\1\{z\in \Gcal\}$ suffices on this channel by~\eqref{eq:family-expfam}; a histogram on bins or $\phi(z)=(1,z,z^2)$ is the general-purpose choice when the channel is unknown.
}

\paragraph{\new{A reference measure as a cost.}}
{\color{blue!70!black}
The reference need not be uniform, and this is where the dependence of~\eqref{eq:ent} on $\mu$ is put to work.
Changing the reference to $\mu'=m\,\mu$ for a positive density $m$ changes the density to $p'=p/m$, and substituting,
\begin{equation}
  H_{\mu'}[p']=-\int\frac{p}{m}\ln\frac{p}{m}\;m\,\dd\mu
  =-\int p\ln p\,\dd\mu+\int p\ln m\,\dd\mu
  =\Hf{p}+\E_{p}\big[\ln m(Z)\big] .
  \label{eq:ref-change}
\end{equation}
A constant $m$ shifts every entropy by the same $\ln m$ and leaves the argmin untouched, which is why the units of $\Zcal$ do not matter; a non-constant $m$ adds a term linear in $\pz$, hence in $\occ$, and does change it.
When $\mu(\Zcal)<\infty$, the choice $m=1/\mu(\Zcal)$ makes $\mu'$ a probability measure, and the entropy of a density against a probability measure is minus its divergence from that measure, so
\begin{equation}
  \Hf{p}=\ln\mu(\Zcal)-\KL\big(p\,\big\|\,\mu/\mu(\Zcal)\big) :
  \label{eq:ent-kl}
\end{equation}
minimising entropy is maximising divergence from the normalised reference, the outcome law that counts as unregulated, uniform on the alphabet for counting $\mu$.
Now take, on a bounded box, $\mu=e^{c}\lambda$ for a bounded function $c:\Zcal\to\mathbb R$, and write $p_\lambda$ for the density of $\pz$ against $\lambda$.
Then~\eqref{eq:ref-change} with $m=e^{c}$ gives
\begin{equation}
  \Hf{\pz}=\Hl{p_\lambda}+\E_{\pz}\big[c(Z)\big],
  \label{eq:ref-cost}
\end{equation}
differential entropy plus an ordinary expected cost, linear in $\occ$.
Minimising~\eqref{eq:obj} under this $\mu$ is therefore regulation plus cost minimisation, and the cost has entered through the yardstick rather than through the channel.
The design rule is that the reference puts mass where outcomes are costly, the opposite of a prior: a Gaussian reference centred on the set-point is $c(z)=-\tfrac12(z-z_0)^2/\sigma^2$ up to a constant and \emph{repels} the regulator from $z_0$, by~\eqref{eq:ent-kl} it maximises the divergence from that Gaussian, whereas $c(z)=+\tfrac12(z-z_0)^2/\sigma^2$ gives entropy plus a quadratic regulation cost.
In the Gibbs form the reward becomes $r_q(s,a)=\E_{K(\cdot\mid s,a)}[\ln q_\lambda(Z)-c(Z)]$, so the cost is an additive reward $-c$ with no iteration-dependent gain, in contrast to $w_k r$ in~\eqref{eq:family-reward}.
The two routes to a reward, through the channel and through the reference, can be combined, and which one a given problem calls for is open.
}

\paragraph{Test setup.}
DeepSea remains the discrete test: an $L\times L$ grid, episodes of exactly $L$ steps, treasure in the bottom-right corner reached only by moving right at every step, with the channel of the first instantiation.
\new{A continuous counterpart with the same structure is a one-dimensional set-point task: the state is a position, actions nudge it, the disturbance is additive noise, and the outcome is emitted by~\eqref{eq:leak-channel} with a band $\Gcal$ that the agent has to find; the box size and $\delta$ set the volume ratio, and the prediction from~\eqref{eq:leak-derivative} is that the learned policy holds the centre of the band with a margin that grows with the noise width.
Both should be run with $\gamma<1$ and with the average-reward estimator side by side, since the theory says they should agree as $\gamma\to1$ on a unichain problem.}

\bibliographystyle{plainnat}
% Overleaf compiles from the project root, a local latexmk -cd from this folder; resolve the .bib paths either way.
\IfFileExists{notes_joery/joery.bib}
  {\newcommand{\bibs}{action-sufficient,notes_joery/joery}}
  {\newcommand{\bibs}{../action-sufficient,joery}}
\bibliography{\bibs}

\end{document}
